\chapter{2005年广东卷}
\section{单选题}
\begin{problem}
若集合 \(M = \{x \mid |x| \leqslant 2\}\)，\(N = \{x \mid x^2 - 3x = 0\}\)，则 \(M \cap N =\)（\quad）

\choices
    {\(\{3\}\)}
    {\(\{0\}\)}
    {\(\{0,2\}\)}
    {\(\{0,3\}\)}
\end{problem}

\begin{answer}
B
\end{answer}

\begin{problem}
若 \((a - 2\i)\i = b - \i\)，其中 \(a\)，\(b \in \R\)，\(\i\) 是虚数单位，则 \(a^2 + b^2 =\)（\quad）

\choices
    {\(0\)}
    {\(2\)}
    {\( \frac{3}{2} \)}
    {\(5\)}
\end{problem}

\begin{answer}
D
\end{answer}

\begin{problem}
\(\displaystyle\lim_{x\to -3}\frac{x + 3}{x^2 - 9} =\)（\quad）

\choices
    {\(-\frac{1}{6}\)}
    {\(0\)}
    {\(\frac{1}{6}\)}
    {\( \frac{1}{3} \)}
\end{problem}

\begin{answer}
A
\end{answer}

\begin{problem}
已知高为 \(3\) 的直棱柱 \(ABC - A'B'C'\) 的底面是边长为 \(1\) 的正三角形（如图所示），则三棱锥 \(B' - ABC\) 的体积为（\quad）

\bitmapfigure[max width=0.34\linewidth]{img/2005/guangdong/q4_fig1.png}

\choices
    {\(\frac{1}{4}\)}
    {\(\frac{1}{2}\)}
    {\(\frac{\sqrt{3}}{6}\)}
    {\(\frac{\sqrt{3}}{4}\)}
\end{problem}

\begin{answer}
D
\end{answer}

\begin{problem}
若焦点在 \(x\) 轴上的椭圆 \(\frac{x^2}{2} + \frac{y^2}{m} = 1\) 的离心率为 \(\frac{1}{2}\)，则 \(m =\)（\quad）

\choices
    {\(\sqrt{3}\)}
    {\(\frac{3}{2}\)}
    {\(\frac{8}{3}\)}
    {\(\frac{2}{3}\)}
\end{problem}

\begin{answer}
B
\end{answer}

\begin{problem}
函数 \( f(x) = x^3 - 3x^2 + 1 \) 是减函数的区间为（\quad）

\choices
    {\((2, +\infty)\)}
    {\((-\infty, 2)\)}
    {\((- \infty, 0)\)}
    {\((0,2)\)}
\end{problem}

\begin{answer}
D
\end{answer}

\begin{problem}
给出下列关于互不相同的直线 \(m\)，\(l\)，\(n\) 和平面 \(\alpha\)，\(\beta\) 的四个命题：

\begin{circlelist}
\item 若 \(m \subset \alpha\)，\(l \cap \alpha = A\)，\(A \notin m\)，则 \(l\) 与 \(m\) 不共面；
\item 若 \(m\)，\(l\) 是异面直线，\(l \parallel \alpha\)，\(m \parallel \alpha\)，且 \(n \perp l\)，\(n \perp m\)，则 \(n \perp \alpha\)；
\item 若 \(l \parallel \alpha\)，\(m \parallel \beta\)，\(\alpha \parallel \beta\)，则 \(l \parallel m\)；
\item 若 \(l \subset \alpha\)，\(m \subset \alpha\)，\(l \cap m = A\)，\(l \parallel \beta\)，\(m \parallel \beta\)，则 \(\alpha \parallel \beta\).
\end{circlelist}
其中为假命题的是（\quad）

\choices
    {\circled{1}}
    {\circled{2}}
    {\circled{3}}
    {\circled{4}}
\end{problem}

\begin{answer}
C
\end{answer}

\begin{problem}
先后抛掷两枚均匀的正方体骰子（它们的六个面分别标有点数 \(1\)，\(2\)，\(3\)，\(4\)，\(5\)，\(6\)），骰子朝上的面的点数分别为 \(X\)，\(Y\)，则 \(\log_{2X}Y = 1\) 的概率为（\quad）

\choices
    {\(\frac{1}{6}\)}
    {\(\frac{5}{36}\)}
    {\(\frac{1}{12}\)}
    {\(\frac{1}{2}\)}
\end{problem}

\begin{answer}
C
\end{answer}

\begin{problem}
在同一平面直角坐标系中，函数 \(y = f(x)\) 和 \(y = g(x)\) 的图象关于直线 \(y = x\) 对称. 现将 \(y = g(x)\) 的图象沿 \(x\) 轴向左平移 2 个单位，再沿 \(y\) 轴向上平移 1 个单位，所得的图象是由两条线段组成的折线（如图所示），则函数 \(f(x)\) 的表达式为（\quad）

\bitmapfigure[max width=0.38\linewidth]{img/2005/guangdong/q9_fig1.png}

\choices
    {\(f(x) = \begin{cases}
2x + 2, & -1\leqslant x\leqslant 0,\\
\frac{x}{2} +2, & 0 < x\leqslant 2.
\end{cases}\)}
    {\(f(x) = \begin{cases}
2x - 2, & -1\leqslant x\leqslant 0,\\
\frac{x}{2} -2, & 0 < x\leqslant 2.
\end{cases}\)}
    {\(f(x) = \begin{cases}
2x - 2, & 1\leqslant x\leqslant 2,\\
\frac{x}{2} +1, & 2 < x\leqslant 4.
\end{cases}\)}
    {\(f(x) = \begin{cases}
2x - 6, & 1\leqslant x\leqslant 2,\\
\frac{x}{2} -3, & 2 < x\leqslant 4.
\end{cases}\)}
\end{problem}

\begin{answer}
A
\end{answer}

\begin{problem}
已知数列 \(\{x_{n}\}\) 满足 \(x_{2} = \frac{x_{1}}{2}\)，\(x_{n} = \frac{1}{2} (x_{n - 1} + x_{n - 2})\)，\(n = 3\)，\(4\)，\(\dots\). 若 \(\displaystyle\lim_{n \to \infty} x_{n} = 2\)，则 \(x_{1} =\)（\quad）

\choices
    {\(\frac{3}{2}\)}
    {\(3\)}
    {\(4\)}
    {\(5\)}
\end{problem}

\begin{answer}
B
\end{answer}

\section{填空题}
\begin{problem}
函数 \(f(x)=\frac{1}{\sqrt{1-\e^x}}\) 的定义域是\fillinblank{}.
\end{problem}

\begin{answer}
\((- \infty,0)\)
\end{answer}

\begin{problem}
已知向量 \(\bs{a} = (2,3)\)，\(\bs{b} = (x,6)\)，且 \(\bs{a} \parallel \bs{b}\)，则 \(x =\fillinblank{}\).
\end{problem}

\begin{answer}
4
\end{answer}

\begin{problem}
已知 \((x\cos \theta + 1)^5\) 的展开式中 \(x^2\) 的系数与 \(\left(x + \frac{5}{4}\right)^4\) 的展开式中 \(x^3\) 的系数相等，则 \(\cos \theta =\fillinblank{}.\)
\end{problem}

\begin{answer}
\(\pm\frac{\sqrt{2}}{2}\)
\end{answer}

\begin{problem}
设平面内有 \(n\) 条直线（\(n \geqslant 3\)），其中有且仅有两条直线互相平行，任意三条直线不过同一点. 若用 \(f(n)\) 表示这 \(n\) 条直线交点的个数，则 \(f(4)=\)\fillinblank{}；当 \(n>4\) 时，\(f(n)=\)\fillinblank{}（用 \(n\) 表示）.
\end{problem}

\begin{answer}
5；\(\frac{n(n-1)}{2}-1\)
\end{answer}

\section{解答题}
\begin{problem}
化简 \(f(x) = \cos\left( \frac{6k + 1}{3} \pi + 2x \right) + \cos\left( \frac{6k - 1}{3} \pi - 2x \right) + 2\sqrt{3} \sin\left( \frac{\pi}{3} + 2x \right)\)（\(x \in \R\)，\(k \in \Z\)），并求函数 \(f(x)\) 的值域和最小正周期.
\end{problem}

\begin{solution}
因为 \(k \in \Z\)，所以
\[
\cos\left(\frac{6k+1}{3}\pi+2x\right)=\cos\left(\frac{\pi}{3}+2x\right)
\qquad
\cos\left(\frac{6k-1}{3}\pi-2x\right)=\cos\left(\frac{\pi}{3}+2x\right).
\]
于是
\[
f(x)=2\cos\left(\frac{\pi}{3}+2x\right)+2\sqrt{3}\sin\left(\frac{\pi}{3}+2x\right)
=4\sin\left(\frac{\pi}{2}+2x\right)=4\cos 2x.
\]
因此 \(f(x)\) 的值域为 \([-4,4]\)，最小正周期为 \(\pi\).
\end{solution}

\begin{problem}
如图所示，在四面体 \(P-ABC\) 中，已知 \(PA = BC = 6\)，\(PC = AB = 10\)，\(AC = 8\)，\(PB = 2\sqrt{34}\). \(F\) 是线段 \(PB\) 上一点，\(CF = \frac{15}{17}\sqrt{34}\)，点 \(E\) 在线段 \(AB\) 上，且 \(EF \perp PB\).

\begin{enumerate}
\item 证明：\(PB \perp\) 平面 \(CEF\)；
\item 求二面角 \(B-CE-F\) 的大小.
\end{enumerate}

\bitmapfigure[max width=0.36\linewidth]{img/2005/guangdong/q16_fig1.png}
\end{problem}

\begin{solution}
建立空间直角坐标系，使 \(B\) 为原点，\(PB\) 为 \(x\) 轴，且 \(A\) 在 \(xOy\) 平面内. 由
\[
PA^2+AB^2=6^2+10^2=136=PB^2
\]
知 \( \triangle PAB\) 在 \(A\) 处为直角三角形. 因此由投影关系，
\[
x_A=\frac{PB^2+AB^2-PA^2}{2PB}=\frac{50}{\sqrt{34}}
\qquad
y_A=\sqrt{AB^2-x_A^2}=\frac{15\sqrt{34}}{17}.
\]
取 \(A\) 在 \(xOy\) 平面的正侧，则
\[
A=\left(\frac{50}{\sqrt{34}},\frac{15\sqrt{34}}{17},0\right).
\]
同理，
\[
x_C=\frac{PB^2+BC^2-PC^2}{2PB}=\frac{18}{\sqrt{34}}.
\]
又由 \(AB^2+BC^2-AC^2=2\overrightarrow{BA}\cdot\overrightarrow{BC}\)，得
\[
\overrightarrow{BA}\cdot\overrightarrow{BC}=\frac{10^2+6^2-8^2}{2}=36.
\]
令 \(C=(x_C,y_C,z_C)\)，则
\[
y_C=\frac{36-x_Ax_C}{y_A}
 =\frac{36-\dfrac{50}{\sqrt{34}}\dfrac{18}{\sqrt{34}}}{\dfrac{15\sqrt{34}}{17}}
 =\frac{27\sqrt{34}}{85},
\]
并且
\[
z_C=\sqrt{BC^2-x_C^2-y_C^2}
 =\sqrt{36-\left(\frac{18}{\sqrt{34}}\right)^2-\left(\frac{27\sqrt{34}}{85}\right)^2}
 =\frac{24}{5}.
\]
取 \(C\) 在 \(xOy\) 平面的正侧，可得
\[
C=\left(\frac{18}{\sqrt{34}},\frac{27\sqrt{34}}{85},\frac{24}{5}\right).
\]

设 \(E=\lambda A\)，其中 \(0\leqslant\lambda\leqslant1\). 因为 \(E\in AB\)，\(F\in PB\)，且 \(EF\perp PB\)，所以
\[
E=\left(\lambda\frac{50}{\sqrt{34}},\lambda\frac{15\sqrt{34}}{17},0\right)
\qquad
F=\left(\lambda\frac{50}{\sqrt{34}},0,0\right).
\]
由 \(CF=\frac{15}{17}\sqrt{34}\)，有
\[
\frac{450}{17}=36-2\lambda\frac{18}{\sqrt{34}}\frac{50}{\sqrt{34}}+\lambda^2\left(\frac{50}{\sqrt{34}}\right)^2
\]
即
\[
1250\lambda^2-900\lambda+162=2(25\lambda-9)^2=0.
\]
故 \(\lambda=\frac{9}{25}\)，从而
\[
x_F=\frac{9}{25}\cdot\frac{50}{\sqrt{34}}=\frac{18}{\sqrt{34}}=x_C.
\]
所以 \(CF\perp PB\). 又已知 \(EF\perp PB\)，而 \(CF\)、\(EF\) 是平面 \(CEF\) 内相交的两条直线，故
\[
PB\perp\text{平面 }CEF.
\]

由 \(\lambda=\frac{9}{25}\)，
\[
E=\left(\frac{18}{\sqrt{34}},\frac{27\sqrt{34}}{85},0\right)
\qquad
F=\left(\frac{18}{\sqrt{34}},0,0\right).
\]
直线 \(CE\) 平行于 \(z\) 轴，因此 \(BE\perp CE\)，\(EF\perp CE\). 设二面角 \(B-CE-F\) 为 \(\theta\)，则其平面角为 \(\angle BEF\)，于是
\[
\cos\theta=\frac{EF}{BE}
 =\frac{\dfrac{27\sqrt{34}}{85}}{\sqrt{\left(\dfrac{18}{\sqrt{34}}\right)^2+\left(\dfrac{27\sqrt{34}}{85}\right)^2}}
 =\frac{\dfrac{27\sqrt{34}}{85}}{\dfrac{306}{85}}
 =\frac{3}{\sqrt{34}}.
\]
所以二面角 \(B-CE-F\) 的大小为 \(\arccos\frac{3}{\sqrt{34}}\).
\end{solution}

\begin{problem}
在平面直角坐标系 \(xOy\) 中，抛物线 \(y = x^2\) 上异于坐标原点 \(O\) 的两不同动点 \(A\)，\(B\) 满足 \(AO \perp BO\)（如图所示）.

\begin{enumerate}
\item 求 \(\triangle AOB\) 的重心 \(G\)（即三角形三条中线的交点）的轨迹方程；
\item \(\triangle AOB\) 的面积是否存在最小值？若存在，请求出最小值；若不存在，请说明理由.
\end{enumerate}

\bitmapfigure[max width=0.34\linewidth]{img/2005/guangdong/q17_fig1.png}
\end{problem}

\begin{solution}
设 \(A=(a,a^2)\)，\(B=(b,b^2)\). 因为 \(A\)，\(B\) 异于 \(O\)，所以 \(a\ne0\)，\(b\ne0\). 由 \(AO\perp BO\)，得
\[
(a,a^2)\cdot(b,b^2)=ab+a^2b^2=ab(1+ab)=0.
\]
因此 \(ab=-1\)，即 \(b=-\frac1a\).

设三角形 \(AOB\) 的重心为 \(G=(u,v)\)，则
\[
3u=a+b=a-\frac1a
\qquad
3v=a^2+b^2=a^2+\frac1{a^2}=\left(a-\frac1a\right)^2+2=9u^2+2.
\]
所以重心 \(G\) 的轨迹方程为
\[
v=3u^2+\frac23.
\]
又因为方程 \(a-\frac1a=3u\) 对任意 \(u\in\R\) 都有非零实根 \(a\)，故轨迹中的 \(u\) 可以取遍 \(\R\). 将坐标变量仍记为 \(x,y\)，轨迹方程为
\[
y=3x^2+\frac23.
\]

三角形 \(AOB\) 的面积为
\[
S_{\triangle AOB}=\frac12\left|\det\begin{pmatrix}a&a^2\\b&b^2\end{pmatrix}\right|
 =\frac12|ab^2-a^2b|
 =\frac12|ab(b-a)|
 =\frac12|a-b|
 =\frac12\left|a+\frac1a\right|.
\]
由
\[
\left(a+\frac1a\right)^2=a^2+2+\frac1{a^2}\geqslant4
\]
得 \(S_{\triangle AOB}\geqslant1\). 当 \(a^2=1\) 时等号成立，此时 \(b=-\frac1a\)，且 \(A\)、\(B\) 确为两不同点. 因此面积存在最小值，最小值为 \(1\).
\end{solution}

\begin{problem}
箱中装有大小相同的黄，白两种颜色的乒乓球，黄，白乒乓球的数量比为 \(s:t\). 现从箱中每次任意取出一个球，若取出的是黄球则结束，若取出的是白球，则将其放回箱中，并继续从箱中任意取出一个球，但取球的次数最多不超过 \(n\) 次，以 \(\xi\) 表示取球结束时已取到白球的次数.
\begin{enumerate}
\item 求 \(\xi\) 的分布列；
\item 求 \(\xi\) 的数学期望.
\end{enumerate}
\end{problem}

\begin{solution}
设
\[
p=P(\text{取到黄球})=\frac{s}{s+t}
\qquad
q=P(\text{取到白球})=\frac{t}{s+t}.
\]
每次取到白球后都放回，因此各次取球相互独立. 若 \(0\leqslant k\leqslant n-1\)，则 \(\xi=k\) 表示前 \(k\) 次取到白球，第 \(k+1\) 次取到黄球；若 \(\xi=n\)，则表示前 \(n\) 次全部取到白球. 因此
\[
P(\xi=k)=
\begin{cases}
q^kp=\left(\dfrac{t}{s+t}\right)^k\dfrac{s}{s+t},&0\leqslant k\leqslant n-1,\\
q^n=\left(\dfrac{t}{s+t}\right)^n,&k=n.
\end{cases}
\]

对取值为非负整数的随机变量使用尾和公式. 对 \(1\leqslant j\leqslant n\)，事件 \(\{\xi\geqslant j\}\) 等价于前 \(j\) 次都取到白球，所以
\[
E(\xi)=\sum_{j=1}^{n}P(\xi\geqslant j)=\sum_{j=1}^{n}q^j
 =\frac{q(1-q^n)}{1-q}
 =\frac{t}{s}\left(1-\left(\frac{t}{s+t}\right)^n\right).
\]
\end{solution}

\begin{problem}
设函数 \(f(x)\) 在 \((-\infty, +\infty)\) 上满足 \(f(2 - x) = f(2 + x)\)，\(f(7 - x) = f(7 + x)\)，且在闭区间 \([0,7]\) 上，只有 \(f(1) = f(3) = 0\).
\begin{enumerate}
\item 试判断函数 \(y = f(x)\) 的奇偶性；
\item 试求方程 \(f(x) = 0\) 在闭区间 \([-2005, 2005]\) 上的根的个数，并证明你的结论.
\end{enumerate}
\end{problem}

\begin{solution}
将题设中的两个对称关系分别改写为
\[
f(x)=f(4-x)
\qquad
f(x)=f(14-x).
\]
于是
\[
f(x)=f(4-x)=f(14-(4-x))=f(x+10)
\]
所以 \(10\) 是函数 \(f(x)\) 的一个周期.

若 \(f(x)\) 是偶函数，则由 \(f(1)=0\) 得 \(f(-1)=0\). 利用周期性，\(f(9)=f(-1)=0\)，再由关于 \(x=7\) 对称的关系得 \(f(5)=f(9)=0\)，这与 \([0,7]\) 上只有 \(1\)，\(3\) 是零点矛盾. 因此 \(f(x)\) 不是偶函数.

若 \(f(x)\) 是奇函数，则 \(f(0)=0\)，从而 \(f(10)=0\). 再由关于 \(x=7\) 对称的关系得 \(f(4)=f(10)=0\)，同样与题设矛盾. 因此 \(f(x)\) 不是奇函数. 故 \(f(x)\) 既不是奇函数，也不是偶函数.

在区间 \([7,10]\) 上，若 \(x\) 是零点，则 \(14-x\in[4,7]\)，且 \(f(x)=f(14-x)\). 由题设，\([4,7]\) 内没有零点，所以 \([0,10)\) 内的零点恰为 \(1\)，\(3\). 结合周期性，方程 \(f(x)=0\) 的全部实根为
\[
x=10k+1\quad\text{或}\quad x=10k+3,\qquad k\in\Z.
\]
在 \([-2005,2005]\) 内，
\[
-2005\leqslant10k+1\leqslant2005
\quad\Longrightarrow\quad -200\leqslant k\leqslant200
\]
以及
\[
-2005\leqslant10k+3\leqslant2005
\quad\Longrightarrow\quad -200\leqslant k\leqslant200.
\]
每类根各有 \(401\) 个，所以根的个数为 \(401+401=802\).
\end{solution}

\begin{problem}
在平面直角坐标系中，已知矩形 \(ABCD\) 的长为 \(2\)，宽为 \(1\)，\(AB\)，\(AD\) 边分别在 \(x\) 轴，\(y\) 轴的正半轴上，\(A\) 点与坐标原点重合（如图所示）. 将矩形折叠，使 \(A\) 点落在线段 \(DC\) 上.

\begin{enumerate}
\item 若折痕所在直线的斜率为 \(k\)，试写出折痕所在直线的方程；
\item 求折痕的长的最大值.
\end{enumerate}

\bitmapfigure[max width=0.34\linewidth]{img/2005/guangdong/q20_fig1.png}
\end{problem}

\begin{solution}
设折叠后 \(A\) 落在线段 \(DC\) 上的点为 \(A'=(r,1)\)，其中 \(0\leqslant r\leqslant2\). 折痕是线段 \(AA'\) 的垂直平分线，因此折痕的斜率为 \(-r\). 若折痕斜率为 \(k\)，则 \(r=-k\)，从而
\[
-2\leqslant k\leqslant0.
\]
折痕经过 \(AA'\) 的中点 \(\left(-\frac{k}{2},\frac12\right)\)，故其方程为
\[
y-\frac12=k\left(x+\frac{k}{2}\right)
\qquad\text{即}\qquad y=kx+\frac{1+k^2}{2}.
\]

令 \(t=-k\)，则 \(0\leqslant t\leqslant2\)，折痕方程为
\[
y=-tx+\frac{1+t^2}{2}.
\]
记折痕在矩形内的长度为 \(L(t)\). 当 \(0\leqslant t\leqslant2-\sqrt3\) 时，折痕与两条竖边相交；当 \(2-\sqrt3\leqslant t\leqslant1\) 时，折痕与左边和下边相交；当 \(1\leqslant t\leqslant2\) 时，折痕与上边和下边相交. 因此
\[
x_0=\frac{1+t^2}{2t}\quad(y=0),\qquad
x_1=\frac{t^2-1}{2t}\quad(y=1),\qquad
y(0)=\frac{1+t^2}{2}.
\]
其中 \(x_0=2\) 的有效解为 \(t=2-\sqrt3\)，\(y(0)=1\) 的有效解为 \(t=1\)，这正是上述三种交点情形的分界.
\[
L(t)=
\begin{cases}
2\sqrt{1+t^2},&0\leqslant t\leqslant2-\sqrt3,\\
\dfrac{(1+t^2)^{3/2}}{2t},&2-\sqrt3\leqslant t\leqslant1,\\
\dfrac{\sqrt{1+t^2}}{t},&1\leqslant t\leqslant2.
\end{cases}
\]
第一段在区间上递增，第三段递减；第二段的导数为
\[
L'(t)=\frac{\sqrt{1+t^2}(2t^2-1)}{2t^2}.
\]
所以第二段在两个端点处取得最大值. 比较端点值，
\[
L(2-\sqrt3)=4\sqrt{2-\sqrt3}>\sqrt2=L(1)
\]
故折痕长的最大值为
\[
4\sqrt{2-\sqrt3}=2\sqrt6-2\sqrt2.
\]
此时 \(t=2-\sqrt3\)，即折痕斜率 \(k=\sqrt3-2\).
\end{solution}
